\textbf{a) 供需平衡。}当段购电、光伏和储能放电应覆盖负荷与充电需求。引入未利用电量$w_t$，将题目中的“不低于”写成等式：
\begin{equation}
 g_t+V_t+b_t=L_t+c_t+w_t,\qquad t\in\mathcal T.\label{eq:balance}
\end{equation}
$w_t\ge0$既可表示无法吸收的光伏，也可表示不接收的已付费外网电量。外网电费仍按$g_t$支付，因此不能再从目标函数中扣掉$p_tw_t$。

\textbf{b) 储电量递推。}充电量和放电量均在交流母线侧计量，输入1\,kWh只能增加0.9\,kWh库存，输出1\,kWh则需要消耗$1/0.9$\,kWh库存：
\begin{equation}
 S_{t+1}=S_t+\eta_c c_t-b_t/\eta_d,\qquad \eta_c=\eta_d=0.9.\label{eq:soc}
\end{equation}
144个区间对应145个状态，$S_t$为本段开始库存，不使用段末值。

\textbf{c) 充放电功率限制。}由于决策变量是电量，5000\,kW的功率上限必须先乘十分钟的时长：
\begin{equation}
 0\le c_t,b_t\le P_{\max}\Delta,\qquad P_{\max}\Delta=5000/6\ \kwh.\label{eq:bounds}
\end{equation}
即使储电量充足，单段放电仍不能超过833.33\,kWh；功率和容量分别约束放电速度与可持续时间。

\textbf{d) 安全库存限制。}在每个区间边界均需检查
\begin{equation}
 1200\le S_t\le10800,\qquad t=0,\ldots,144.\label{eq:inventory}
\end{equation}
额定容量为12000\,kWh，但实际可以调度的库存幅度为9600\,kWh，不能将额定容量全部用于套利。

\textbf{e) 日循环边界。}给定初始库存为6000\,kWh，问题一每天重复同一组输入，因此要求
\begin{equation}
 S_0=S_{144}=6000.\label{eq:cycle}
\end{equation}
仅给日末库存设置物理下限，会允许消耗期初储能来降低当日费用，不能代表可重复运行的日计划。

\textbf{f) 购电方向及运行模式。}微网不向外网售电，余量不能为负，同一段储能只允许一种运行方向：
\begin{equation}
 g_t,w_t\ge0,\qquad c_tb_t=0.\label{eq:nonnegative}
\end{equation}
前两个不等式限制交易和余电的方向，乘积约束限制设备状态，不能用购电非负代替充放电互斥。
